\subsubsection{P5：灵敏度分析}

P5 对应运筹学中的\textbf{灵敏度分析、容量规划和决策支持}。在 P1--P4 已形成可行调度链路的基础上，P5 进一步研究关键参数变化对目标函数 $C_{\max}$、总延期和可行性的影响。由于 AMR 数量 $m$、任务量 $n$ 和瓶颈容量 $q_C$ 均为离散变量，P5 采用\textbf{情景重优化 + 有限差分 + 方差分解}的方法，对边际收益和主导影响因素进行定量评估。

\paragraph{离散资源边际分析。}
P5 将 AMR 数量视为整数资源变量，对不同 $m$ 逐点重新求解调度问题，并用前向差分衡量增加 1 台 AMR 对系统性能的边际改善。这对应运筹学中资源配置变量的边际分析。若 $\Delta C_{\max}$ 随 $m$ 增加而快速减小，说明继续增加 AMR 的边际收益递减，系统逐渐进入资源相对充足状态；若改善幅度仍然明显，则说明 AMR 数量仍是限制系统性能的重要资源约束。

\paragraph{系统容量边界与可行性判定。}
任务负荷分析对应\textbf{容量规划}问题。P5 在同一任务池上构造嵌套前缀，并对不同任务规模逐点重优化；同时结合最晚完成时刻、延期任务数、路径缺失和电量违规构造可行性判据，再通过 all-seeds 规则确定保守容量边界。该过程体现了\textbf{资源受限调度}和\textbf{可行域判定}思想，即先判定系统在给定任务负荷下能否稳定完成服务，再进一步评价服务质量和性能水平。

\paragraph{瓶颈容量与影子价值。}
对局部通道簇容量 $q_C$ 的放松实验，对应\textbf{离散影子价格分析}。当某一瓶颈通道的容量增加 1 个单位后，重新执行冲突修复与调度求解，进而观察 $C_{\max}$ 和总等待时间的下降量。该下降量可用于衡量局部通道容量的边际价值，并为瓶颈识别、通道扩容和资源优先配置提供依据。

\paragraph{Morris 与 Sobol 全局灵敏度。}
P5 的全局筛选使用 Morris 和 Sobol 两类经典灵敏度分析方法。Morris 方法通过基本效应 $\mu^\star$ 和波动度 $\sigma$ 筛选主导因素，适用于离散、非线性和具有局部交互的系统；Sobol 方法通过方差分解计算一阶效应和总效应，用于区分单因素主效应与多因素交互效应。两者结合后，可以同时识别关键参数及其交互影响，从而提高 P5 对调度系统性能变化来源的解释能力。
